% !PRINTED 260
\weberpage{260}

A number of the form $a\varepsilon$ shall be called a \gesperrt{purely imaginary}
\gesperrt{number}, and two numbers of the form $a+b\varepsilon$, $a-b\varepsilon$ \gesperrt{con-}
\gesperrt{jugate imaginary numbers}. The transfer of these designations
from the theory of the ordinary complex numbers
to the numbers of the field $\mathfrak{E}_{2,p}$ is justified by the
agreement of the rules of calculation, and can lead to no misunderstanding,
since in these considerations there is no
question of the ordinary complex numbers.

Since always, when $N$ is a quadratic non-residue,
$N^{1/2(p-1)}=-1$, it follows from (2):

\weq{3}{\varepsilon^p = -\varepsilon.}

If one applies the binomial theorem to the $p$th power of the
binomial $\alpha=a+b\varepsilon$, and observes that all the binomial coefficients,
with the exception of the first and the last, are divisible
by $p$, and hence here are $=0$, and further that by \gesperrt{Fermat's}
theorem $a^p=a$, $b^p=b$, then by (3) we obtain:

\weq{4}{\alpha = a+b\varepsilon, \quad \alpha^p = a-b\varepsilon,}
and consequently, by multiplication:

\weq{5}{\alpha^{p+1} = a^2 - Nb^2,}
from which it follows that $\alpha^{p+1}$ is always a real number.

\gesperrt{If $\gamma$ is a primitive root of the field $\mathfrak{E}_{2,p}$ [\S.\,64, (8)],
then $\gamma^r$ is real when and only when $r$ is divisible by $p+1$,
and $\gamma^{p+1}$ is a primitive root of the
prime $p$ in the ordinary sense}.

\gesperrt{The necessary and sufficient condition
for a number $\alpha$ in $\mathfrak{E}_{2,p}$ to be real consists in the
equation $\alpha^p=\alpha$}.

\gesperrt{Every real number is representable as the $(p+1)$th power of a number
in $\mathfrak{E}_{2,p}$}.

For that $\gamma^r$ is real when $r$ is divisible by $p+1$
follows from (5). But that it can be real in no other way results
from the fact that by Fermat's theorem every real number
different from zero satisfies the condition $a^{p-1}=1$, so that,
if $\gamma^r$ is real, $\gamma^{r(p-1)}=1$ must hold, and therefore
$r(p-1)$ must be divisible by $p^2-1$, and consequently $r$ by $p+1$.
Since, further, $\gamma^{r(p+1)}$ is $=1$ only when $r$ is divisible by $p-1$,
$\gamma^{p+1}$ is a primitive root of $p$.

\gesperrt{The concept of the roots of unity can also be transferred to
the numbers of the field $\mathfrak{E}_{2,p}$}.
